\section{Anketna pitanja i oznake odgovora}

\subsection{Pitanje P9}

\textbf{P9:} Što smatrate najvećom prednošću SMR tehnologije?

Ponuđeni odgovori:

\begin{itemize}[nosep]
  \item Niski početni i ukupni troškovi investicije
  \item Daleko kraće vrijeme izgradnje i puštanje u pogon
  \item Poboljšane sigurnosne značajke reaktora
  \item Ne znam
\end{itemize}

\subsection{Pitanje P10}

\textbf{P10:} Smatrate li opravdanim državnu subvenciju razvoja SMR tehnologije umjesto ulaganja u jednu klasičnu, veliku nuklearnu elektranu?

Ponuđeni odgovori:

\begin{itemize}[nosep]
  \item Da
  \item Ne
  \item Ne znam
\end{itemize}

\subsection{Legenda skraćenih oznaka}

\begin{table}[H]
\centering
\caption{Legenda skraćenih oznaka korištenih u grafikonima}
\begin{tabular}{ll}
\toprule
Oznaka & Puni odgovor \\
\midrule
Niži troškovi & Niski početni i ukupni troškovi investicije \\
Kraća izgradnja & Daleko kraće vrijeme izgradnje i puštanje u pogon \\
Sigurnosne značajke & Poboljšane sigurnosne značajke reaktora \\
Ne znam & Ne znam \\
\bottomrule
\end{tabular}

\end{table}

\newpage

\section{Pseudokod statističke analize}

U ovom dodatku prikazan je sažeti pseudokod stvarne statističke analize provedene nad anketnim podacima. Izostavljeni su tehnički dijelovi programa koji se odnose samo na spremanje datoteka, oblikovanje grafova i pripremu izlaza za izvještaj.

Analiza se temelji na opisnoj statistici i hi-kvadrat testovima nezavisnosti. Opisna statistika korištena je za prikaz raspodjele odgovora, dok su hi-kvadrat testovi korišteni za provjeru postoji li statistički značajna povezanost između demografskih obilježja i odgovora na pitanja P9 i P10 te između samih odgovora na pitanja P9 i P10.

\begin{enumerate}[label=\textbf{Korak \arabic*:}, leftmargin=*]
  \item Učitati anketne podatke iz Excel datoteke.
  \item Odabrati varijable korištene u analizi:
  \begin{itemize}[nosep]
    \item demografske varijable: županija, spol, dob, obrazovanje i zaposlenost,
    \item odgovore na pitanja P9 i P10.
  \end{itemize}
  \item Očistiti tekstualne odgovore uklanjanjem suvišnih razmaka i praznih vrijednosti.
  \item Za pitanja P9 i P10 zamijeniti duge tekstualne odgovore kraćim oznakama radi preglednijeg prikaza u tablicama i grafovima.
  \item Za svaku analiziranu varijablu izračunati frekvencijsku tablicu:
  \[
    \text{postotak kategorije} = \frac{\text{broj odgovora u kategoriji}}{\text{ukupan broj valjanih odgovora}} \cdot 100.
  \]
  \item Definirati parove varijabli za testiranje povezanosti:
  \begin{itemize}[nosep]
    \item svaka demografska varijabla s pitanjem P9,
    \item svaka demografska varijabla s pitanjem P10,
    \item pitanje P9 s pitanjem P10.
  \end{itemize}
  \item Za svaki par varijabli ukloniti zapise u kojima nedostaje odgovor na barem jednu od te dvije varijable.
  \item Prije testiranja spojiti kategorije koje imaju manje od pet opažanja u zajedničku kategoriju \textit{Ostalo (<5)}. Time se smanjuje rizik nepouzdanih rezultata zbog vrlo malih očekivanih frekvencija.
  \item Izraditi kontingencijsku tablicu za promatrani par varijabli.
  \item Ako barem jedna varijabla nakon čišćenja ima manje od dvije kategorije, preskočiti hi-kvadrat test za taj par jer test nije smislen.
  \item Inače provesti hi-kvadrat test nezavisnosti nad kontingencijskom tablicom:
  \[
    \chi^2 = \sum \frac{(O - E)^2}{E},
  \]
  gdje je $O$ opažena frekvencija, a $E$ očekivana frekvencija uz pretpostavku nezavisnosti varijabli.
  \item Za svaki provedeni test zabilježiti vrijednost $\chi^2$, broj stupnjeva slobode, $p$-vrijednost i najmanju očekivanu frekvenciju.
  \item Izračunati Cramérov $V$ kao mjeru jačine povezanosti:
  \[
    V = \sqrt{\frac{\chi^2}{n \cdot \min(r - 1, k - 1)}},
  \]
  gdje je $n$ ukupan broj opažanja u kontingencijskoj tablici, $r$ broj redaka, a $k$ broj stupaca.
  \item Parove varijabli smatrati statistički značajno povezanima ako je $p < 0{,}05$.
  \item Za statistički značajne parove izraditi grafički prikaz postotaka po retku kako bi se vidjelo kako se raspodjela odgovora mijenja između kategorija prve varijable.
\end{enumerate}

Rezultati dobiveni ovim postupkom korišteni su u glavnom dijelu izvještaja za tablice frekvencija, tablicu hi-kvadrat testova i prikaze statistički značajnih povezanosti.
